\section{Discussion}

\label{sec:discussion}

\subsection{What the hybrid is good for}

The result is better read as a successful audit: in the canonical setup, 38 of
    6{,}680 decisions change, each described by a channel, a direction and a
mechanism.

\emph{Triage and accountability.} A queue can sort comments by whether the decision
was rule-supported: a \hate{} decision with a matched span and $r_H=0.8$ is easier
to review than one with no lexical evidence, and the span, the entry and the
arithmetic are what a platform can defend in policy terms.

\emph{Editing and diagnosis.} Adding a slur variant is a configuration change, not
a retraining run, and an edit affects the \hate{} channel (weight 0.8) far more than
the others. The decomposition exists only because the channels are separate: a
monolithic classifier improving by 0.0027 macro-F1 would give no handle on whether
the change was lexical or neural.

\subsection{Why accuracy does not move}

\label{sec:why-neutral}

First, \emph{coverage is low}: the rule fires on 15.75\% of comments, so a
perfect rule channel could still only change decisions in that subset. Second,
\emph{the majority class dominates}: with 83\% \clean{}, only the 1{,}132 toxic
comments can improve, and 38 changes is 3.4\% of that support---enough to
matter operationally, too little to move a macro-average beyond its confidence
interval. Third, \emph{tuning is safety-oriented}: protecting \clean{} and
gating \hate{} on lexical support produces both the de-flagging and the missed
corrections. Neutrality is a \emph{choice}, not an absence of signal. The
leak-free sensitivity run shows a similar decision pattern, with 36 changes (23
corrections and 13 regressions) and macro-F1 0.6612; it is reported as a
separate sensitivity result rather than merged with the canonical analysis.

\subsection{The directional asymmetry is a policy question}

\label{sec:asymmetry}

In a binary ``remove / keep'' framing the asymmetry of
Sect.~\ref{sec:regressions} is a recall-oriented trade-off; in a three-class
framing it is not, because an \off{} comment mislabelled as \hate{} can trigger
a stronger sanction than the content warrants, and a political comment
mislabelled as \hate{} is a different kind of harm.

That false positive matters: the entry is a charged multi-word expression
labelled \hate{} from co-occurrence, so the lexicon is not a neutral
artefact---it encodes editorial judgments that enter at weight 0.8. Deployments
should separate \emph{profanity} from \emph{political} entries and weight them
accordingly.

\subsection{Design implications}

\label{sec:implications}

Four changes follow. (i) \textbf{Evidence-gated fusion}: when the rule fires
strongly against a high-confidence \clean{} prediction, abstain or escalate.
(ii) \textbf{Separate thresholds per mechanism}: requiring $r_H\ge0.5$ before
the \hate{} channel can override would make the rule-driven and
threshold-driven paths separately tunable. (iii) \textbf{Per-class calibration}
of the \clean{} channel, which currently blocks the missed corrections. (iv)
\textbf{Human-in-the-loop} for the 804 both-wrong comments, which no fusion of
these two channels can fix.

\subsection{Limitations}

\label{sec:limitations}

\paragraph{Single benchmark and language.}
All results come from ViHSD and Vietnamese. The lexicon and cue taxonomy are
language-specific; the fusion protocol is not, but we have not tested it
elsewhere.

\paragraph{Preliminary three-seed check.}
The main backbone sweep was run with seed 42. A preliminary three-seed check of
the focused fusion comparison produced macro-F1 differences (fusion minus
AI-only) of $+0.0027$, $-0.0041$ and $-0.0010$, for a mean difference of
$-0.0008$. The sign changes across runs, and the mean is slightly negative; the
single-seed gain is therefore not stable across this preliminary check. These
runs do not replace a full multi-seed evaluation of every backbone and fusion
variant.

\paragraph{Cross-corpus overlap.}
The exact-text audit found 895 of 6{,}680 ViHSD test rows (13.4\%) with matches
in the auxiliary corpus material examined for lexicon development; 620 distinct
test texts also occur in the Bontori/R-ViHSD training partition. This overlap
creates a potential source of lexical leakage for the canonical lexicon.
Rebuilding the lexicon after excluding the overlapping Bontori/R-ViHSD training
texts produced a 671-entry leak-free lexicon. The resulting standalone
rule-engine and fusion scores are reported separately in
Section~\ref{sec:leak-free-sensitivity}. This sensitivity analysis addresses
the identified overlap but does not establish the absence of all possible
cross-corpus dependencies.

\paragraph{Repeated test texts.}
The frozen test split contains 6{,}680 rows but 6{,}571 distinct texts; one
text appears twice among the 15 canonical regressions. We retain all rows to
preserve the fixed split, but no deduplicated sensitivity analysis or
duplicate-group bootstrap was available. The reported row-level resampling
therefore may not fully account for dependence among repeated texts.

\paragraph{The rule engine is weaker than designed.}
Validation set the accent-less alias weight to zero and the negation factor to
one, removing the main defence against dropped diacritics and the discount for
negated slurs visible in the \off{}-as-intensifier regressions.

\paragraph{Explanations are not evaluated by humans.}
The evidence record is faithful for the lexical channel and decomposes the
changed decisions, but we do not measure whether reviewers find it useful, and
the cue taxonomy is a post-hoc, unblinded annotation.

\subsection{Ethics}

\label{sec:ethics}

The comments are public ViHSD test comments and are masked as described in
Sect.~\ref{sec:data}. The system is decision support, not automatic
sanctioning: the directional asymmetry and the political false positive show
that automatic \hate{} labels from this pipeline should be reviewed before they
affect a user. The lexicon should be treated as a governed artefact with a
change log, because its entries shape moderation outcomes.